\documentclass[12pt]{article}
\usepackage{graphicx}
\usepackage{hyperref}
\usepackage{enumitem}

\begin{document}

Ilmenau Technical University

Institute for Media and Communication Science

Predicting Student Dropout in Higher Education Using Machine Learning and Digital Engagement Data

A Systematic Literature Review Grounded in Learning Analytics

By: Moritz Tesel, Ehrenbergstr. 29, 98693 Ilmenau

Date: 09.01.2026


\tableofcontents

\newpage

\section{Introduction}

Student dropout represents a persistent challenge for higher education institutions worldwide, with substantial individual and societal costs. The increasing digitalization of higher education has produced rich digital engagement data that can support early identification of students at risk.

\section{State of Research}

Early research on student dropout prediction relied on demographic and academic background variables. More recent studies show a clear shift toward behavioural data from learning management systems.

\section{Theoretical Framework}

This study conceptualizes student dropout as the outcome of a longitudinal process of disengagement. Learning analytics provides the conceptual foundation for measuring that process.

\section{Methodology}

\subsection{Study Design}

The study adopts a systematic literature review design following established guidelines for evidence synthesis in educational research.

\begin{table}[h]
\caption{PICOC framework for the systematic literature review}
\label{tab:picoc}
\centering
\begin{tabular}{ll}
Population & Students enrolled in higher education programmes \\
Intervention & Machine learning models using digital engagement data \\
Comparison & Traditional statistical baselines \\
Outcome & Prediction of dropout or retention \\
Context & Higher education institutions \\
\end{tabular}
\end{table}

The review is guided by the following research questions:

\begin{enumerate}[label=RQ\arabic*:]
\item What types of data sources and digital engagement features are used in machine learning based dropout prediction?
\item Which machine learning algorithms are most commonly applied?
\item How is prediction performance evaluated, and which evaluation metrics are most frequently reported?
\end{enumerate}

\subsection{Search Strategy and Selection Criteria}

The search covered three databases and yielded 625 records. After duplicate removal and screening, 42 studies met the inclusion criteria.

\begin{figure}[h]
\centering
\includegraphics[width=0.8\textwidth]{prisma_diagram}
\caption{PRISMA 2020 Flow Diagram of the Study Selection Process}
\label{fig:prisma}
\end{figure}

\section{Results}

\subsection{Data Sources and Digital Engagement Features}

Learning management system logs are the most common data source, followed by student information systems.

\subsection{Machine Learning Algorithms}

Random forests and logistic regression are the most frequently applied algorithms, with neural approaches gaining ground.

\subsection{Evaluation Metrics and Performance}

Model quality is most frequently reported using accuracy, precision, recall, and the F1-score\footnote{The F1-score is defined as follows: \(F_1 = 2 \cdot \frac{\mathrm{Precision} \cdot \mathrm{Recall}}{\mathrm{Precision} + \mathrm{Recall}}\). Its highest possible value is 1.0, indicating perfect precision and recall.}, complemented by the ROC-AUC\footnote{The Receiver Operating Characteristics (ROC) curve plots the True Positive Rate against the False Positive Rate for all possible classification thresholds.}.

\section{Discussion}

The findings confirm the central role of behavioural engagement data while exposing weaknesses in evaluation practice.

\section{Conclusion}

Machine learning based dropout prediction has matured, but methodological rigour varies widely across studies.

\section{References}

Tinto, V. (1975). Dropout from Higher Education. Review of Educational Research.

Romero, C. and Ventura, S. (2020). Educational Data Mining and Learning Analytics. WIREs Data Mining and Knowledge Discovery.

\end{document}
